\par\bigskip\Needspace{8\baselineskip}
\pdfbookmark[1]{Störungen der Ceres: Bemerkungen}{ceres-bemerkungen}
\SourceLeaf{416}
\begin{center}BEMERKUNGEN.\end{center}
\begingroup\small
\noindent\hspace*{\parindent}Die Notiz [2] des Abschnitts I. ist einem Heftchen Schedae, die Notizen [3]---[7] einem Handbuche entnommen; wegen der fragmentarischen Beschaffenheit der Originale musste die Bearbeitung ziemlich frei gehalten werden. Das genannte Handbuch enthält auch die erste numerische Rechnung der Störungen nach den in [7] gegebenen Formeln, deren Resultate in \textsc{Zachs} Monatlicher Correspondenz Bd. VI (Werke, Bd. VI S. 225) abgedruckt sind. Auf Grund dieser Störungen hat \textsc{Gauss} die Elemente wiederum verbessert (Bd. VI S. 228 und Notiz [1]), und dann mit den so verbesserten Elementen die Störungen von neuem berechnet; diese letztere Rechnung findet sich in einigen zusammengehefteten Blättern, welche den Titel »Praktische Anweisung zur Berechnung der Störungen« tragen, und Eingangs auch eine Zusammenstellung der angewandten Formeln mit Einschluss derer für die Breitenstörungen (Notiz [8]) enthalten. Die letztern hatte \textsc{Gauss} in der ersten Rechnung bei Seite gelassen; ihre Ableitung hat sich nicht gefunden. Die Resultate der zweiten Rechnung sind abgedruckt in \textsc{Zachs} Monatlicher Correspondenz, Bd. VI (Werke Bd. VI S. 228---229). Die Jupitersmasse hat \textsc{Gauss} gleich \(\dfrac{1}{1067{,}09}\) (nach \textsc{Laplace}) angenommen.

Die unter II. abgedruckten Briefstellen beziehen sich auf die Tafeln, welche \textsc{Gauss} im Anschluss an die vorerwähnten Rechnungen hergestellt und in \textsc{Zachs} Monatlicher Correspondenz Bd. VII (Werke Bd. VI S. 235---243) veröffentlicht hat.

In der zwischen den beiden Briefen an \textsc{Olbers} vom 25. Januar und 10. Mai 1805 liegenden Zeit (oder früher) hat \textsc{Gauss} begonnen, seine erste Methode auf die zweiten Potenzen der Excentricitäten auszudehnen. Der Beginn dieser Entwickelungen findet sich ebenfalls im genannten Handbuch; sie sind aber ihrer Weitläuftigkeit wegen von \textsc{Gauss} sehr bald wieder aufgegeben worden, worauf dann die Rechnung nach der zweiten unter III. mitgetheilten Methode erfolgte. Vgl. hierzu den Brief an \textsc{Olbers} vom 10. Mai 1805.

Von dieser zweiten Methode, bei welcher \textsc{Gauss} bereits die spätere \textsc{Hansen}sche interpolatorische Entwickelung der Störungsfunction angewandt hat (vgl. \textsc{Gauss} an \textsc{Hansen}, 11. März 1843), und die auch die Grundlage der spätern Rechnung der Pallasstörungen bildet, sind nur die numerischen Rechnungen in demselben Handbuch aufzufinden gewesen, auf Grund deren die Entwickelungen reconstruirt worden sind. Da demnach die Notizen [2]---[5] gänzlich neu verfasst sind, so sind sie in Petitsatz gedruckt. Aus dem Originale stammen nur: S. 402 die Werthe von \(N\), \(\sin J\), \(\AscendingNode\), und die Erklärung der Bedeutung von \(\alpha,\epsilon,\epsilon'\); S. 403 die Gleichung 2) und der Werth der Constante bei Gleichung 3); S. 405---407 die numerischen Werthe mit Ausnahme derer auf S. 405, Zeile 9---5 von unten, und der Factoren von \(\cos(3D+\text{\AstroSymbols ⚳}-\AscendingNode+M)\) auf S. 406, Zeile 11---15, welche ergänzt wurden. \textsc{Gauss} hat hier nur die Breitenstörungen gerechnet; von der Anwendung auf die Störungen des Radius Vector und der Länge fand sich nur der allererste Anfang vor, \textsc{Gauss} ist dann zur Bearbeitung der Pallasstörungen übergegangen. Die Resultate finden sich im Briefe an \textsc{Olbers} vom 2. Juli 1805; sie sind damals nicht abgedruckt worden.

Im ganzen haben \textsc{Gauss}’ Arbeiten über Ceres den Charakter von Vorarbeiten zu seinen ausgedehnten Untersuchungen über Pallas, womit es sich wohl rechtfertigt, dass manche Punkte dort ausführlicher bearbeitet worden sind, als hier.

Man vergleiche auch die bereits in Bd. VI abgedruckten einzelnen Mittheilungen über den Fortgang von \textsc{Gauss}’ Untersuchungen. Mehrfach wurden kleinere Unrichtigkeiten verbessert.
\begin{flushright}\textsc{Brendel}.\end{flushright}
\par\endgroup
\begin{center}\rule{.38\linewidth}{.4pt}\end{center}
